\section{Check if There Exists a Subsequence with Sum K}
\label{sec:rec-check-subsequence-sum-k}

\problemstatement{Given an array $\textit{arr}$ and an integer $k$,
determine whether \textbf{any} subsequence of $\textit{arr}$ sums to
exactly $k$.}

\begin{patternbox}
This is Section~\ref{sec:rec-power-set}'s pick/not-pick tree, but the
question changes from \emph{``generate every subsequence''} to
\emph{``does at least one satisfy a condition.''} Whenever a backtracking
problem only needs a yes/no answer (rather than every solution), the
recursion should \textbf{stop exploring the instant an answer is found}
-- there is no need to generate the other $2^n-1$ subsequences once one
of them has already confirmed the answer is ``yes.''
\end{patternbox}

\subsection*{Understanding the problem carefully}

Instead of building and recording a list at every leaf, track a running
sum as the pick/not-pick decisions are made, and check it only once
every element has been decided (at the base case). The function should
return a boolean, and -- critically -- the moment either branch reports
\texttt{true}, the calling level should return \texttt{true} immediately
without bothering to explore the other branch at all.

\begin{ideabox}[Same Tree, but Short-Circuit on Success]
$\textsc{Exists}(\textit{index}, \textit{sum})$: if
$\textit{index} = n$: return $(\textit{sum} = k)$. Otherwise: if
$\textsc{Exists}(\textit{index}+1, \textit{sum} +
\textit{arr}[\textit{index}])$ (the pick branch) is \texttt{true}, return
\texttt{true} immediately without checking the not-pick branch at all.
Otherwise, return $\textsc{Exists}(\textit{index}+1, \textit{sum})$ (the
not-pick branch) directly as this call's result.
\end{ideabox}

\subsection*{Why short-circuiting never misses a valid answer}

We only need to know whether \emph{some} subsequence achieves sum $k$ --
not which one, nor how many do. The moment the pick branch reports
success, we already have a definitive ``yes'' for the overall question,
and continuing to explore the not-pick branch (or any other untried
branch) could only ever confirm the same ``yes'' again, never change it
to a ``no.'' Skipping that redundant exploration is therefore always
safe, and in the best case it can save an enormous amount of work,
since a successful subsequence might be found very early in the
traversal while the vast majority of the $2^n$ possibilities are never
even visited.

\subsection*{Algorithm}

\begin{enumerate}
  \item $\textsc{Exists}(\textit{index}, \textit{sum})$:
  \begin{enumerate}
    \item If $\textit{index} = n$: return $(\textit{sum} = k)$.
    \item If $\textsc{Exists}(\textit{index}+1, \textit{sum} +
          \textit{arr}[\textit{index}])$: return \texttt{true}.
    \item Return $\textsc{Exists}(\textit{index}+1, \textit{sum})$.
  \end{enumerate}
  \item Call $\textsc{Exists}(0, 0)$.
\end{enumerate}

\subsection*{Dry run}

\dryrun{Take $\textit{arr} = [1,2,3]$, $k=5$.}

\begin{itemize}
  \item $\textsc{Exists}(0,0)$: pick $1$: $\textsc{Exists}(1,1)$.
  \item $\textsc{Exists}(1,1)$: pick $2$: $\textsc{Exists}(2,3)$.
  \item $\textsc{Exists}(2,3)$: pick $3$: $\textsc{Exists}(3,6)$: base
        case, $6 \ne 5$, return \texttt{false}. Not-pick $3$:
        $\textsc{Exists}(3,3)$: base case, $3 \ne 5$, return
        \texttt{false}. So $\textsc{Exists}(2,3)$ returns
        \texttt{false}.
  \item Back in $\textsc{Exists}(1,1)$: pick branch was \texttt{false},
        try not-pick $2$: $\textsc{Exists}(2,1)$: pick $3$:
        $\textsc{Exists}(3,4)$: base, $4 \ne 5$, false. Not-pick $3$:
        $\textsc{Exists}(3,1)$: base, $1\ne5$, false. So
        $\textsc{Exists}(2,1)$ is false, and $\textsc{Exists}(1,1)$ is
        false.
  \item Back in $\textsc{Exists}(0,0)$: pick branch was \texttt{false},
        try not-pick $1$: $\textsc{Exists}(1,0)$: pick $2$:
        $\textsc{Exists}(2,2)$: pick $3$: $\textsc{Exists}(3,5)$: base,
        $5=5$! Return \texttt{true}. This propagates straight back up:
        $\textsc{Exists}(2,2)$ returns \texttt{true} immediately (no
        need to check not-pick $3$); $\textsc{Exists}(1,0)$ returns
        \texttt{true} immediately; $\textsc{Exists}(0,0)$ returns
        \texttt{true}.
\end{itemize}

The answer is \texttt{true}: the subsequence $\{2,3\}$ sums to $5$, and
the recursion stopped exploring further branches the instant it was
found.

\subsection*{C++ implementation}

\begin{lstlisting}
#include <bits/stdc++.h>
using namespace std;

bool exists(int index, int sum, vector<int>& arr, int k) {
    if (index == (int)arr.size()) return sum == k;

    if (exists(index + 1, sum + arr[index], arr, k)) return true;
    return exists(index + 1, sum, arr, k);
}

int main() {
    vector<int> arr = {1, 2, 3};
    cout << (exists(0, 0, arr, 5) ? "true" : "false") << endl;
    return 0;
}
\end{lstlisting}

\begin{complexitybox}
\textbf{Time Complexity.} $O(2^n)$ in the worst case (no qualifying
subsequence exists, so every branch must be explored), but often much
faster in practice due to short-circuiting on the first success.
\textbf{Space Complexity.} $O(n)$ for the recursion stack.
\end{complexitybox}

\begin{takeawaybox}
``Does there exist'' questions should always short-circuit the moment a
valid answer is found -- propagate a \texttt{true} result immediately up
through every level of the recursion, skipping any sibling branch that
has not yet been explored. This is a small but important variation on
the pick/not-pick template: the same tree shape, but with an early-exit
discipline that the ``generate everything'' version of
Section~\ref{sec:rec-power-set} does not need.
\end{takeawaybox}
